\section*{Chapitre \ref{ch:g2:born-grow-die} --- Naître, grandir, mourir}

\begin{solution}{exo:g2:born-grow-die:1}
Naître, être jeune, être \omterm{def:g2:born-grow-die:cycle}{adulte}, vieillir (et à la fin, mourir).
\end{solution}

\begin{solution}{exo:g2:born-grow-die:2}
À l'\omterm{def:g2:born-grow-die:cycle}{étape} \omterm{def:g2:born-grow-die:cycle}{adulte}.
\end{solution}

\begin{solution}{exo:g2:born-grow-die:3}
Du plus âgé au plus jeune~: la grand-mère, la mère, l'écolier, le
bébé.
\end{solution}

\begin{solution}{exo:g2:born-grow-die:4}
Un noyau de cerise germe (la \omterm{def:g2:born-grow-die:cycle}{naissance})~; le jeune arbre grandit
pendant des années~; l'arbre \omterm{def:g2:born-grow-die:cycle}{adulte} donne des \omterm{def:g1:plants-around-us:parts}{fleurs} et des cerises
--- donc de nouveaux noyaux~; puis il vieillit et meurt.
\end{solution}

\begin{solution}{exo:g2:born-grow-die:5}
Par exemple~: un bébé a des dents, les perd et en a de plus
solides~; le duvet jaune du poussin devient de vraies \omterm{def:g1:animals-around-us:covering}{plumes}~; la
\omterm{def:g1:plants-around-us:parts}{tige} verte et tendre d'un jeune arbre devient du bois dur.
\end{solution}

\begin{solution}{exo:g2:born-grow-die:6}
De la plus courte à la plus longue~: le papillon (quelques
semaines), le chat (dix à vingt ans), le chêne (des centaines
d'années).
\end{solution}

\begin{solution}{exo:g2:born-grow-die:7}
Parce qu'avant la fin de la route, les \omterm{def:g2:born-grow-die:cycle}{adultes} ont des petits, et
que les petits recommencent la route --- c'est la flèche violette
qui reboucle. Une vie est une route~; la vie au fil des générations
tourne comme une roue.
\end{solution}

\begin{solution}{exo:g2:born-grow-die:8}
Réponse libre. Les photos montrent d'habitude le bébé, l'enfant,
l'\omterm{def:g2:born-grow-die:cycle}{adulte} --- la \omterm{def:g2:born-grow-die:cycle}{vieillesse} est peut-être l'\omterm{def:g2:born-grow-die:cycle}{étape} encore à venir.
\end{solution}

\begin{solution}{exo:g2:born-grow-die:9}
Les pommes du vieil arbre portent des pépins --- des graines --- et
chaque pépin planté est un jeune pommier~: la génération suivante.
Quand le vieil arbre meurt, ses petits continuent le cycle, et le
jardin garde ses pommiers.
\end{solution}

\begin{solution}{exo:g2:born-grow-die:10}
Non --- c'est parce que le caillou n'a jamais été vivant. Seuls les
\omterm{def:g1:living-or-not:living}{êtres vivants} naissent, grandissent et meurent~; un caillou n'a pas
de \omterm{def:g2:born-grow-die:cycle}{cycle de vie} du tout, donc «~ne jamais mourir~» n'est pas de la
solidité, c'est ne pas \omterm{def:g1:living-or-not:living}{être vivant}.
\end{solution}
